\clearpage
\section{附录：Homework 源码与说明}

本附录收录课程作业 2--6 的完整源码与设计说明。

% ===== Homework 2 =====
\subsection{Homework 2：并行归并排序}

\noindent\textbf{题目：}长度为 $m$ 的数组二分 $k$ 次，得到 $2^k$ 个子数组，每个子数组内部排序后两两归并。进程数 $n = \min(2^k, m)$，$m$ 和 $k$ 运行时输入。

\lstset{basicstyle=\scriptsize\fontspec{DejaVu Sans Mono},language=C++}
\lstinputlisting{../homework2/homework2.cpp}

\noindent\textbf{运行示例：}
\lstset{basicstyle=\scriptsize\fontspec{DejaVu Sans Mono},language={}}
\begin{lstlisting}
mpiexec -n 4 homework2.exe
m=8, k=2, needed processes=4, actual=4
Rank 0 got chunk: 3 1
Rank 0 sorted: 1 3
Rank 0 merged from rank 1, now has: 1 1 3 4
Rank 0 merged from rank 2, now has: 1 1 2 3 4 5 6 9
Final sorted array: 1 1 2 3 4 5 6 9
PASS: matches std::sort
\end{lstlisting}

\clearpage
% ===== Homework 3 =====
\subsection{Homework 3：MPI 矩阵乘法}

\noindent\textbf{题目：}利用 9 个进程并行计算 $3\times3$ 矩阵乘法 $C = A \times B$，每个进程算一个元素。

\lstset{basicstyle=\scriptsize\fontspec{DejaVu Sans Mono},language=C++}
\lstinputlisting{../homework3/homework3.cpp}

\noindent\textbf{运行示例：}
\lstset{basicstyle=\scriptsize\fontspec{DejaVu Sans Mono},language={}}
\begin{lstlisting}
mpiexec -n 9 homework3.exe
Matrix A:
   1   2   3
   4   5   6
   7   8   9
Matrix B:
   9   8   7
   6   5   4
   3   2   1
C = A x B:
  30  24  18
  84  69  54
 138 114  90
\end{lstlisting}

\clearpage
% ===== Homework 4 =====
\subsection{Homework 4：用 MPI\_Send / MPI\_Recv 实现 MPI\_Gather}

\noindent\textbf{题目：}不使用 MPI\_Gather，仅用 MPI\_Send / MPI\_Recv 实现自定义 \texttt{my\_MPI\_Gather}，并与标准 MPI\_Gather 对比验证。

\lstset{basicstyle=\scriptsize\fontspec{DejaVu Sans Mono},language=C++}
\lstinputlisting{../homework4/homework4.cpp}

\noindent\textbf{设计思路：}
\begin{itemize}
  \item root 进程先将自身数据拷贝到 recvbuf 头部。
  \item 然后按 rank 顺序接收其他进程发送的数据，写入 recvbuf 对应偏移位置。
  \item 非 root 进程各发送一个 MPI\_INT 到 root。
  \item 用标准 MPI\_Gather 的结果做校验。
\end{itemize}

\clearpage
% ===== Homework 5 =====
\subsection{Homework 5：用 MPI\_Send / MPI\_Recv 实现 MPI\_Allgather}

\noindent\textbf{题目：}不使用 MPI\_Allgather，仅用 MPI\_Send / MPI\_Recv 实现自定义 \texttt{my\_MPI\_Allgather}，并与标准 MPI\_Allgather 对比验证。

\lstset{basicstyle=\scriptsize\fontspec{DejaVu Sans Mono},language=C++}
\lstinputlisting{../homework5/homework5.cpp}

\noindent\textbf{设计思路（环形轮转）：}
\begin{enumerate}
  \item 每个进程先将自己的数据块写入 recvbuf 对应位置（rank 索引处）。
  \item 进行 $size-1$ 步环形传递：当前进程向 (rank+1) 发送一个数据块，同时从 (rank-1) 接收一个数据块。
  \item 每步目标地址由 (rank - step) mod size 计算，保证 $size-1$ 步后所有数据到达所有进程。
\end{enumerate}

\clearpage
% ===== Homework 6 =====
\subsection{Homework 6：思政思考——华为在高性能/并行计算领域的核心竞争力}

\lstset{basicstyle=\scriptsize\fontspec{DejaVu Sans Mono},language={}}
以下为思政作业原文：

\begin{lstlisting}
# 思政思考：华为在高性能/并行计算领域的核心竞争力

## 与并行计算课程的联系

华为在 HPC/并行计算领域最大的竞争力在于
"芯片—框架—编译器"全栈自研的并行计算生态，与课程内容直接对应：

| 课程概念                  | 华为对应能力                                     |
|--------------------------|------------------------------------------------|
| MPI 集体通信（Gather、Allgather、Reduce 等） | MindSpore 框架内置的分布式并行策略，自动在集群中完成梯度 Allreduce、参数 Gather 等操作 |
| 进程间通信、拓扑结构      | HCCS（华为高速互联）实现昇腾芯片间低延迟通信，类似 MPI 中 communicator 的物理层实现 |
| 并行效率、负载均衡        | 昇腾 Ascend 的达芬奇架构（DaVinci Core）专为矩阵运算并行设计，3D Cube 单元单次可完成 16x16x16 的矩阵乘加 |
| 编译器优化                | 毕昇编译器（Bisheng）针对鲲鹏/昇腾自动向量化、循环并行化 |

## 华为独特壁垒

相比 MPI 只定义"通信标准"的开放模式，华为的独特竞争力在于软硬协同：

- 鲲鹏处理器提供 ARM 架构的通用算力
- 昇腾 AI 处理器提供专用并行加速（达芬奇架构的 3D Cube）
- MindSpore 自动并行策略编排，开发者无需手动写 MPI 通信代码
- 毕昇编译器针对华为硬件深度优化

三者全栈打通，从芯片到并行框架全面自主可控。
\end{lstlisting}
